%%%% LISTA DE ABREVIATURAS E SIGLAS 
%%
%% Relação, em ordem alfabética, das abreviaturas (representação de uma palavra por meio de alguma(s) de sua(s) sílaba(s) ou
%% letra(s)), siglas (conjunto de letras iniciais dos vocábulos e/ou números que representa um determinado nome) e acrônimos
%% (conjunto de letras iniciais dos vocábulos e/ou números que representa um determinado nome, formando uma palavra pronunciável).
%%
%% Este arquivo para definição de abreviaturas, siglas e acrônimos é utilizado com a opção "glossaries" (pacote)

%% Abreviaturas: \abreviatura{rótulo}{representação}{definição}

% \listadeabrevsiglaseacr

% % \abreviatura{art.}{art.}{Artigo}
% % \abreviatura{cap.}{cap.}{Capítulo}
% % \abreviatura{sec.}{sec.}{Seção}

% %% Siglas: \sigla{rótulo}{representação}{definição}

% \sigla{abnt}{ABNT}{Associação Brasileira de Normas Técnicas}
% \sigla{cnpq}{CNPq}{Conselho Nacional de Desenvolvimento Científico e Tecnológico}
% \sigla{eps}{EPS}{\textit{Encapsulated PostScript}}
% \sigla{pdf}{PDF}{Formato de Documento Portátil, do inglês \textit{Portable Document Format}}
% \sigla{ps}{PS}{\textit{PostScript}}
% \sigla{utfpr}{UTFPR}{Universidade Tecnológica Federal do Paraná}


%% LEIA:

%% Para usar o \gls, você deve colocar a sigla aqui em \sigla

%% ADICIONAR SIGLAS: Quando você inclui alguma sigla, pode ser necessário compilar umas duas vezes para essa aparecer na lista de siglas.

%% ATENÇÃO REMOVER SIGLAS: se você remover a sigla do seu texto (não for usar mais), você deve comentar essa aqui e remover os \gls{} dessa sigla (se não vai ficar aparecendo a sigla na lista). Em caso de ERRO, quando o LaTeX informa que você ainda tem a sigla no texto, mesmo que não tenha. Você deve limpar o cache - no OverLeaf, clique no erro, vá:
%  ->view error
%%   ->(role para baixo, até o final)
%%     ->e clique em Clear cached files


%% Siglas do trabalho

%% BEAM é definida sem expansão no corpo do texto. O nome por extenso é um
%% acrônimo retroativo, longo e pouco informativo no meio da frase; a forma
%% completa continua registrada na Lista de Abreviaturas e Siglas.
%% Por isso a entrada é declarada aqui em vez de usar \siglaIt: o campo "first"
%% recebe a sigla, e não a expansão.
\newglossaryentry{BEAM}{
 type=\pseudoacronymtype,
 name=BEAM,
 description={\textit{Bogdan/Björn's Erlang Abstract Machine}},
 first={BEAM},
 long={\textit{Bogdan/Björn's Erlang Abstract Machine}}
}
\siglaPt{VM}{Máquina Virtual}
%% NIF é declarada aqui, e não com \siglaIt, para ter formas de plural próprias.
%% Com a definição genérica, "\gls{NIF}s" saía como "Native Implemented Function
%% (NIF)s", com o "s" fora dos parênteses. Use \glspl{NIF} no plural.
\newglossaryentry{NIF}{
 type=\pseudoacronymtype,
 name=NIF,
 plural={NIFs},
 description={\textit{Native Implemented Function}},
 first={\textit{Native Implemented Function} (NIF)},
 firstplural={\textit{Native Implemented Functions} (NIFs)},
 long={\textit{Native Implemented Function}}
}
\siglaIt{Nx}{Numerical Elixir}
\siglaIt{EXLA}{Elixir XLA}
\siglaIt{XLA}{Accelerated Linear Algebra}
\siglaIt{API}{Application Programming Interface}
\siglaIt{GIL}{Global Interpreter Lock}
%% RSS é definida sem expansão no corpo do texto: o significado já é dado em
%% português ("uso de memória residente") nos pontos em que aparece, e a
%% expansão em inglês nesses trechos produziria parênteses aninhados.
\newglossaryentry{RSS}{
 type=\pseudoacronymtype,
 name=RSS,
 description={\textit{Resident Set Size}},
 first={RSS},
 long={\textit{Resident Set Size}}
}
\siglaIt{OTP}{Open Telecom Platform}
\siglaIt{ML}{Machine Learning}
\siglaIt{IPC}{Inter-Process Communication}

%% Acrônimos: \acronimo{rótulo}{representação}{definição}
%\acronimo{gimp}{Gimp}{Programa de Manipulação de Imagem GNU, do inglês \textit{GNU Image Manipulation Program}}
